\documentclass[10pt,a4paper]{amsart}
\usepackage[margin=19mm]{geometry}
\usepackage{amsmath,amssymb,mathtools}
\usepackage[scr=rsfs,cal=euler]{mathalfa}
\usepackage{xfrac}
\usepackage[unicode,hidelinks,bookmarks=false]{hyperref}
\newcommand{\ie}{i.e.\ }
\newcommand{\cf}{cf.\ }
\newcommand{\resp}{respectively\ }
\newcommand{\Subsection}{\subsection}
\providecommand{\nobreakdash}{\nobreak\hbox{-}}
\setlength{\emergencystretch}{2em}
\multlinegap0pt
\usepackage{bm}
\usepackage{bbm}
\usepackage{dsfont}
\usepackage{xspace}
\usepackage{cancel}
\usepackage{mathtools}
\usepackage{amsthm}
\makeatletter
\@ifundefined{theorem}{\newtheorem{theorem}{Theorem}}{}
\@ifundefined{lemma}{\newtheorem{lemma}[theorem]{Lemma}}{}
\makeatother
\newcommand{\EE}{\mathbb{E}} % 
\newcommand{\Ew}{\mathbf{E}_{\omega}} % quenched law
\newcommand{\NN}{\mathbb{N}}
\newcommand{\Pws}{\mathrm{P}_{\omega}} % quenched law of RW
\newcommand{\TT}{\mathbb{T}} % Ulam tree
\newcommand{\ind}[1]{\mathbbm{1}_{\left\{ #1 \right\} } } %indicator function
\newcommand{\rt}{\mathbf{o}} % the root of the GW tree
\title[Biased branching random walks on Bienaym\'e--Galton--Watson ]{Biased branching random walks on Bienaym\'e--Galton--Watson trees}
\author{Julien Berestycki, Nina Gantert, David Geldbach, Quan Shi}
\date{Source: arXiv:2502.07363v3 (CC BY 4.0)}
\begin{document}
\maketitle

\noindent\emph{Source and licence.} Biased branching random walks on Bienaym\'e--Galton--Watson trees. Version arXiv:2502.07363v3 (CC BY 4.0), distributed under the Creative Commons Attribution 4.0 International licence, as recorded in the arXiv licence metadata for this exact version and shown on the abstract page https://arxiv.org/abs/2502.07363v3. Full rights notice: licenses/arxiv-2502.07363-CC-BY-4.0.txt.\par\smallskip

\noindent\textit{Editorial scope.} Counted unit: the Lemma whose source label is \texttt{lem:log} in the article (source0.tex lines 1059--1070, source label \texttt{lem:log}), delivered together with the article's own complete original proof of it (source0.tex lines 1072--1100, one \texttt{proof} environment of 29 lines). No mathematics is written, repaired or re-derived: statement and proof are the article's own text line for line. Required context retained uncounted: eq:many\_to\_one (lines 581--583); eq:Z1 (lines 895--897); lem:log0-bis (lines 1104--1111). Every definition, display and label of those lines is retained unchanged. Omitted from this package and not redistributed: 1--580, 584--894, 898--1058, 1071--1071, 1101--1103, 1112--1338. Nothing inside a retained excerpt is omitted or altered: every retained body is delivered exactly as the article's lines are written, so no omission receipt of any kind accompanies this package. Citations: the retained excerpts carry no citation command at all, so no bibliography entry of the article is transcribed and no reference is redistributed. Reference adaptation: none was needed and none was made; every reference inside the delivered unit resolves inside it, so no unresolved reference or citation is produced. Printed slips kept literal: no printed word, symbol, hypothesis or number of the counted statement or of its proof was altered. Layout and class: the article's own class is \texttt{article} with packages including amssymb, amsmath, amsfonts, amsthm, geometry, mathrsfs, xcolor, url, enumitem, bbm, lineno, mathtools, hyperref, titlesec; the wrapper is the lane's pinned \texttt{amsart} editorial wrapper at 10pt on A4, which restates the theorem-like environments the retained text needs and adds the article's own macro definitions that the retained text needs (\texttt{\textbackslash EE}, \texttt{\textbackslash Ew}, \texttt{\textbackslash NN}, \texttt{\textbackslash Pws}, \texttt{\textbackslash TT}, \texttt{\textbackslash ind}, \texttt{\textbackslash rt}), with extra wrapper packages (bm, bbm, dsfont, xspace, cancel, mathtools). The delivered numbering is the unit's own. Assets: no retained span contains a figure environment, a graphics-inclusion command, a PDF-page inclusion, a table or a TikZ picture; no image file is copied into this package, the package declares no asset dependency and no image file is redistributed with it. Licence: version arXiv:2502.07363v3 of this article is distributed under the Creative Commons Attribution 4.0 International licence, as recorded in the arXiv licence metadata for this exact version and shown on the abstract page https://arxiv.org/abs/2502.07363v3. A full rights notice is delivered at licenses/arxiv-2502.07363-CC-BY-4.0.txt. Characters: every retained excerpt is pure ASCII, so the delivered engine, XeLaTeX with its default Latin Modern text font, reports no missing character.
\begin{equation}\label{eq:many_to_one}
	\mathbf{E}_{\omega}\left[\sum_{u\in \mathbb{T}\colon\vert u \vert = n} f\left((X(u_k), k\leq n) \right)\right] = m^n \mathrm{E}_{\omega} \left[ f\left((S_k, k\leq n) \right) \right], \qquad n\in \NN.
\end{equation}

\begin{equation}\label{eq:Z1}
\mathcal{Z}^{(n)}_{1} = \{u\in  \mathcal{L}_n \colon |u| \le a^{-1} n,X(u_k)\ne \rt, \forall k= 1,\ldots, |u|   \} . 
\end{equation}

\begin{lemma}\label{lem:log0-bis}
	Let $(S_j, j\ge 0)$ be the $\lambda$--biased random walk starting from $\rt$. Set $T_n = \inf\{j\ge 0\colon |S_j|=n  \}$ and $\tau_{\rt} = \inf\{j > 0\colon S_j = \rt  \}$. 
	Let $a,b\in (v_{\lambda}, 1]$ with $a<b$. 
	Then we have 
	\begin{equation*}
		\lim_{n\to \infty} \EE_{\mathtt{BGW}}  \Big[ \frac{a}{n}\log  \Pws\big( b^{-1} n< T_{n} \le a^{-1}  n, \tau_\rt >T_{n}  \big)  \Big]  = - I_{\lambda}(a). 
	\end{equation*}
\end{lemma}

\begin{lemma}\label{lem:log}
  	Let  $v_{\lambda} < a< v_{\lambda, m}$. 
  		Recall from \eqref{eq:Z1} that $\mathcal{Z}^{(n)}_1:= \{u\in \mathcal{L}_n \colon |u| \le a^{-1} n,X(u_k)\ne \rt, \forall k= 1,\ldots, |u| \} $. 
  	It holds that
	\begin{equation}
	\liminf_{n\to \infty}\EE_{\mathtt{BGW}} \left[\frac{a}{n}\log \Ew\big[\# \mathcal{Z}^{(n)}_1 \big] \right] >0.
	\end{equation}
	In particular, there is $n^* \in \NN$ such that for all $n\geq n^*$ we have
	\begin{equation}\label{eq:lem:log}
	\EE_{\mathtt{BGW}} \left[\log \Ew\big[	\# \mathcal{Z}^{(n)}_1   \big] \right] >0. 
	\end{equation}
\end{lemma}

\begin{proof}
	For $v\in \TT$ with $\vert v\vert\ge n$, we write $\tau_\rt^v > n$ if $\forall 0< k\leq n$  
	we have $X(v_k)\neq \rt$. Similarly, for the $\lambda$--biased random walk $(S_j, j\geq 0)$ we set $T_n = \inf\{j\ge 0\colon |S_j|=n  \}$ and $\tau_{\rt} = \inf\{j > 0\colon S_j = \rt  \}$.  
	
	Fix any $\varepsilon>0$. By the many-to-one formula \eqref{eq:many_to_one},
	\begin{align*}
		  & \Ew\left[	\# \{v\in \TT\colon |v|= \lceil a^{-1} n\rceil , (a+\varepsilon)^{-1}n<T^v_{n } \le a^{-1} n  , \tau^v_\rt >T^v_{n}  \} \right] \\
		& = m^{ \lceil a^{-1}  n\rceil } \Pws\big( (a+\varepsilon)^{-1}n< T_{n } \le a^{-1} n, \tau_\rt >T_{n}  \big). 
	\end{align*} 
On the other hand, by the branching property at the stopping line 
$\mathcal{L}_{n}=\{u_{T^u_{n }}, u\in \TT \}$, we have 
\begin{align*}
	&\Ew\left[\# \{v\in \TT\colon  |v|= \lceil a^{-1} n\rceil ,(a+\varepsilon)^{-1}n<T^v_{n } \le a^{-1} n , \tau^v_\rt >T^v_{n}  \} \right] \\
	&= \Ew\left[	\sum_{u\in \mathcal{L}_n } m^{\lceil a^{-1} n \rceil - |u|} \ind{u\in\mathcal{Z}^{(n)}_1 , |u|>  (a+\varepsilon)^{-1}n} \right] 
		\leq m^{\lceil a^{-1}n \rceil- (a+\varepsilon)^{-1}n}\Ew\big[	\# \mathcal{Z}^{(n)}_1 \big].
\end{align*}
Combining the two observations, taking the logarithm and rearranging leads to
	\begin{align}\label{eq:lem-log-Z1}
	&	\log\Ew\big[	\# \mathcal{Z}^{(n)}_1\big]
		\ge (a+\varepsilon)^{-1} n \log m + \log  \Pws\big( (a+\varepsilon)^{-1}n< T_{n } \le a^{-1} n, \tau_\rt >T_{n}  \big) . 
	\end{align}
 We integrate both sides in $\omega$ with respect to $\mathtt{BGW}$ and let $n\to \infty$.  
It follows from Lemma \ref{lem:log0-bis} below  that
\begin{equation}
	\liminf_{n \to \infty}\EE_{\mathtt{BGW}} \left[ \frac{a}{n} \log \Ew\big[\# \mathcal{Z}^{(n)}_1\big]\right] \ge \frac{a}{a+\varepsilon} \log m -I_{\lambda} (a).\label{eq:rw-ldp}
\end{equation}
Since $I_{\lambda}$ is strictly increasing on $(v_{\lambda},1]$ with $\log m \leq I_{\lambda}(v_{\lambda, m})$, we have  $\log m> I_{\lambda}(a)$, thus the right-hand side is strictly positive for $\varepsilon$ small enough. 
This completes the proof.
\end{proof}

\end{document}
